\documentclass{article}

\title{\Huge Global Trade Visualizer}
\date{Software Engineering \\  Semester 6 \\  27-02-2020}
\author{Andrei Cioanca 266105 \\ \\ Claudiu Rediu 266129 \\ \\ Stefan Harabagiu 266116 \\ \\  Supervisors: Poul Erik Væggemose\\ \\ VIA University College}
\usepackage{fancyhdr} % Customizable Headers
\usepackage{hyperref}
\usepackage[toc,page]{appendix}
\pagestyle{fancy}
\usepackage{graphicx}
\fancyhf{}
\rhead{\includegraphics[width=75px]{VIAUniversityLogo.pdf}}
\lhead{Global Trade Visualizer}
\rfoot{Page \thepage}
\begin{document}
	\pagenumbering{gobble}
	\maketitle
	\newpage
	\pagenumbering{arabic}
	\tableofcontents
	\newpage
	\section{Project Description}
	Global Trade Visualizer is an application that will make use of readily available online data related to trades between countries. Part of what it will be able to achieve, is providing an environment in which users can visualize and explore different trading behaviours and patterns. This product would be able to switch between a broad and narrow context, so different users with varying expertise can obtain insights.\newline
	It will be able to make use of datasets from providers such as the World Integrated Trade Solution(WITS) that keeps track of traded goods and services between different countries and continents. This data will be processed and it will be part of two flows. The first one will provide for the visualization and the second one will be the support for different reports that will go in depth.\newline
	The product will be able to help in planning and optimising supply chains, together with the opportunity to explore new untouched markets.
	\section{Technologies Description}
	The technologies that are going to be used are:
	\begin{enumerate}
		\item SQL Server for storing the trade data
		\item C\# for building a web service around the app
		\item Python for obtaining and processing the data
		\item Unreal Engine for building the visualization
		\item C++ in Unreal Engine
	\end{enumerate}
	\section{Project Owner}
	Andrei Cioanca will be the product owner that will aid in managing expectations of the customer and the results of the project.
	\section{Project Group}
	The group will be composed of three members. \newline
	Andrei Cioanca has the interactive media specialization.  He has extensive experience in building visualization software using third party 3D game engines such as Unreal Engine as well as developing communications between different applications.\newline
	Claudiu Rediu has the data specialisation. He has experience working close to the supply chain of a company, helping them with data flows that feed the reports and doing technical analysis to optimize procedures.\newline
	Stefan Harabagiu has the data specialisation. He has experience with developing and deploying web services in a fast-paced environment in which it is vital to continuously maintain.
\end{document}